% Source: https://artofproblemsolving.com/community/c3h1367358_plane_geometry_problem
% Stable identity: aops-topic-1367358-q01
Since $CA=CB=CD=10$, the points $A$, $B$, $D$ lie on the circle with centre $C$ and radius $10$. Triangle $ABC$ is equilateral, so the central angle $\angle ACB=60^\circ$, i.e.\ the minor arc $AB$ measures $60^\circ$.

\begin{asy}
usepackage("amsmath");
size(7cm);
defaultpen(fontsize(10pt));
pair C=(0,0); pair A=10*dir(200), B=10*dir(260);
real t=2*aSin(0.65); pair D=10*dir(200-t);
draw(circle(C,10),gray+dashed);
draw(A--B--C--cycle,linewidth(0.9)); draw(C--D,linewidth(0.9));
draw(A--D,blue); draw(B--D,blue);
draw(arc(D,2.6,degrees(B-D),degrees(A-D)),red);
label("$30^\circ$",D+5.0*dir((degrees(B-D)+degrees(A-D))/2),red);
draw(arc(C,1.3,200,260)); label("$60^\circ$",C+2.4*dir(225));
dot(C);
label("$A$",A,W); label("$B$",B,S); label("$C$",C,NE); label("$D$",D,NW);
label("$13$",(A+D)/2,W,blue); label("$10$",(C+D)/2,NE);
label("\mbox{}",(min(currentpicture,user=true)+max(currentpicture,user=true))/2,black); // keeps the text after the figure black when LaTeX typesets the labels inline
\end{asy}

Where is $D$? The chord $AD=13$ subtends a central angle $\angle ACD=2\arcsin\dfrac{13}{20}\approx81^\circ>60^\circ$. Any point of the minor arc $AB$ is within $60^\circ$ of $A$, so $D$ is not on the minor arc $AB$ (on either side of $A$); it lies on the major arc. By the inscribed angle theorem,
\[\angle ADB=\dfrac{1}{2}\angle ACB=\boxed{30^\circ}.\]
